% wave13_a3_gkm_presentation_etingof.tex
%
% Etingof-channelled adversarial assault + rigorous healing of Lorgat 2020
% Section 5: "The generalized Kac-Moody algebra g and its automorphic
% correction g_{Delta_5}".
%
% FIVE ATTACK <-> HEAL cycles, each mathematically substantive.
%
% Primary sources audited:
%   - Lorgat 2020 (arXiv:2005.05017 source excerpt, wave13_source_paper.txt)
%   - Gritsenko-Nikulin 1995 (arXiv:alg-geom/9504006), Part II
%   - Borcherds 1998, Invent. Math. 132 (singular theta lift, Thm 13.3)
%   - Kac, "Infinite-dimensional Lie algebras" Ch.11 (GKM axioms)
%   - Ray 2007, "Automorphic forms and Lie superalgebras" (BKM superalgebras)
%
% Cross-refs: wave12_u1_two_stage_functor.tex, wave12_a7_coha_w_infty_drinfeld.tex
%             (CoHA = Y^+, not W_{1+infty}), wave12_a2_k3xE_BKM_kazhdan.tex.
%
% Scope: chain-level, object-level g_{Delta_5}. The (infty,1)-categorical
% lane (g_{Delta_5} as CoHA image) is referenced but not re-proved here.

\documentclass[11pt]{article}
\usepackage[localtheorems]{raeez-math-template}

\newtheorem{theorem}{Theorem}
\newtheorem{proposition}[theorem]{Proposition}
\newtheorem{lemma}[theorem]{Lemma}
\newtheorem{corollary}[theorem]{Corollary}
\theoremstyle{definition}
\newtheorem{definition}[theorem]{Definition}
\newtheorem{remark}[theorem]{Remark}
\newtheorem{attack}[theorem]{Attack}
\newtheorem{heal}[theorem]{Heal}

\providecommand{\C}{\mathbb{C}}
\providecommand{\Z}{\mathbb{Z}}
\providecommand{\Q}{\mathbb{Q}}
\providecommand{\R}{\mathbb{R}}
\providecommand{\N}{\mathbb{N}}
\providecommand{\g}{\mathfrak{g}}
\providecommand{\Aut}{\mathrm{Aut}}
\providecommand{\Sp}{\mathrm{Sp}}
\providecommand{\mult}{\mathrm{mult}}
\providecommand{\sdim}{\mathrm{sdim}}
\providecommand{\ad}{\mathrm{ad}}
\providecommand{\Lp}{\Lambda^{2,1}}
\providecommand{\LII}{\Lambda^{2,1}_{I\!I}}
\providecommand{\PII}{\mathcal{P}_{I\!I}}
\providecommand{\CII}{\mathcal{C}(\LII)_+}
\providecommand{\dre}{\Delta^{\mathrm{re}}}
\providecommand{\dim}{\Delta^{\mathrm{im}}}
\providecommand{\kBKM}{\kappa_{\mathrm{BKM}}}
\providecommand{\kcat}{\kappa_{\mathrm{cat}}}
\providecommand{\kch}{\kappa_{\mathrm{ch}}}
\providecommand{\Deltafive}{\Delta_5}
\providecommand{\phionezero}{\phi_{0,1}}

\title{Etingof-voiced adversarial assault and rigorous healing of\\
       Lorgat 2020 \S 5: the GKM superalgebra $\g_{\Deltafive}$}
\author{Wave 13 DNA strand A3}
\date{2026-04-22}

\begin{document}
\maketitle

\begin{abstract}
We audit \S 5 of Lorgat 2020 from first principles in the voice of Pavel
Etingof. Five ATTACK$\leftrightarrow$HEAL cycles verify (i) the rank-$3$
hyperbolic Cartan matrix of signature $(2,1)$; (ii) the Weyl vector
$\rho=\tfrac12(\delta_1+\delta_2+\delta_3)$ and the simple-real-root
equation $(\rho,\alpha)=-\alpha^2/2$; (iii) the two-family stratification
of imaginary simples $\dim=\dim_0\sqcup\dim_1$ (even Gritsenko--Nikulin
null orbits of multiplicity $\tau(a)=9$; odd timelike orbits of
multiplicity $|m(a)|$); (iv) the Macdonald-type identity applied to the
trivial module $\C$ reproducing exactly Gritsenko's infinite product for
$\tfrac{1}{64}\Deltafive(2Z)$; (v) the super-structure: the algebra has
no real odd simples, $\tau(a)=9$ even null simples per Aut$(\PII)=S_3$
orbit, and odd simples supplied by timelike $m(a)<0$ coefficients of
$\phionezero$. We heal the presentation into explicit
Chevalley--Serre--Borcherds-Ray form with generators, relations, and
Weyl--Kac--Borcherds denominator.
\end{abstract}

\section*{First-principles protocol}

\paragraph{(a) The question.} What is the precise GKM-superalgebra
structure on $\g_{\Deltafive}$, and does its Weyl--Kac--Borcherds
denominator identity reconstruct the Gritsenko--Nikulin product
$\tfrac{1}{64}\Deltafive(2Z)$?

\paragraph{(b) Primary sources.}
Gritsenko--Nikulin 1995, arXiv:alg-geom/9504006 (``Siegel automorphic
form corrections of some Lorentzian Kac--Moody Lie algebras''), Part II
\S 2--\S 3, for the $\Deltafive$ root system; Borcherds 1998, Invent.
132, Thm 13.3, for the singular theta lift of weight
$c_{\phionezero}(0,0)/2=10/2=5$; Kac, ``Infinite-dimensional Lie
algebras'' Ch.~11, for the GKM axiom list; Ray 2007 (J. Algebra 307),
for the BKM-superalgebra extension allowing odd imaginary simples.

\paragraph{(c) Chain-level or $(\infty,1)$?}
Chain level throughout. Every coefficient
$f(n,l)$, every multiplicity $\mult\alpha$, every product exponent is
expanded explicitly. The $(\infty,1)$-categorical identification
$\g_{\Deltafive}\simeq$ BKM-image of a $d=3$ CY datum through $\Phi$ is
a distinct statement (cross-ref wave12\_u1\_two\_stage\_functor) and is
not re-proved here.

%=====================================================================
\section{Lattice and Cartan data for $\g$ before the BKM correction}
%=====================================================================

Let $\Lp$ be the hyperbolic lattice of signature $(2,1)$ with basis
$f_2, f_3, f_{-2}$; $\LII\subset\Lp$ is the even sublattice of index $8$.
Let $\PII$ denote the fundamental polyhedron of the even-reflection group
$W^{(2)}(\LII)$, with primitive orthogonal vectors
$\PII^{\mathrm{prim}}=\{\delta_1,\delta_2,\delta_3\}$,
$\delta_1=f_2-f_3$, $\delta_2=f_{-2}-f_2$, $\delta_3=f_3$.
Write $G_{ij}=(\delta_i,\delta_j)$ for the Gram matrix.

\begin{attack}[Cartan signature]
\label{att:cartan-signature}
The paper asserts the Gram matrix
\[
G=(\delta_i,\delta_j)=
\begin{pmatrix} 2 & -2 & -2\\ -2 & 2 & -2\\ -2 & -2 & 2\end{pmatrix}
\]
and silently treats $G$ as the ``symmetric generalized Cartan matrix''
of a rank-$3$ hyperbolic Kac--Moody algebra. Verify:
(a) the off-diagonal Cartan integers
$\tfrac{2(\delta_i,\delta_j)}{(\delta_i,\delta_i)}$ are
$-2\in\Z_{\le 0}$, consistent with Kac GCM axiom (C2);
(b) $G$ has signature $(2,1)$, so $\g(G)$ is a Lorentzian
Kac--Moody algebra in the sense of Gritsenko--Nikulin;
(c) $\det G$ is nonzero, ruling out affine type.
\end{attack}

\begin{heal}[Diagonalisation, signature, hyperbolicity]
\label{heal:cartan-signature}
The vector $v_0=(1,1,1)^\top$ satisfies $G v_0=(-2,-2,-2)^\top=-2\,v_0$,
giving eigenvalue $\lambda_0=-2$. Orthogonal complement in $\R^3$ is the
plane $\{(x,y,z):x+y+z=0\}$. For $v_1=(1,-1,0)^\top$:
$Gv_1=(4,-4,0)^\top=4v_1$. For $v_2=(1,1,-2)^\top$:
$Gv_2=(4,4,-8)^\top=4v_2$. Hence $\mathrm{Spec}(G)=\{-2,4,4\}$, signature
$(+,+,-)=(2,1)$, $\det G = -32$. The algebra $\g$ associated via Kac Ch.1
to $G$ (equivalently: to the generalised Cartan matrix
$A_{ij}=2G_{ij}/(\delta_i,\delta_i)$, which here equals $G$ because
$(\delta_i,\delta_i)=2$) is the rank-$3$ hyperbolic Kac--Moody algebra
of Lorentzian signature.

This $\g$ is already infinite-dimensional; its root system has real
simple roots $\dre=\{\delta_1,\delta_2,\delta_3\}$ but \emph{no
imaginary simples}. The Kac imaginary cone
$\{\alpha:(\alpha,\alpha)\le 0,\ \alpha\in\CII\}$ contains
countably-infinitely many roots all of the same Kac multiplicity; their
total contribution to the $\g$-denominator cannot reconstruct
$\Deltafive$ because $\Deltafive$ has $\Sp_4(\Z)$-automorphy which the
$\g$-denominator lacks. This gap --- real and quantitative --- is what
forces the Borcherds--Ray automorphic correction to $\g_{\Deltafive}$.
\end{heal}

%=====================================================================
\section{Weyl vector and the simple-real-root equation}
%=====================================================================

\begin{attack}[Weyl vector existence]
\label{att:weyl-vector}
The paper writes a lattice Weyl vector
$\rho\in\PII\otimes\Q$ satisfying $(\rho,\delta_i)=-1$ for
$i=1,2,3$. Derive this from first principles: does the linear system
$(\rho,\delta_i)=-\tfrac{(\delta_i,\delta_i)}{2}=-1$ have a unique
solution in $\LII\otimes\Q$, and is that solution in the positive cone
$\CII$? Without $\rho\in\CII$ the Weyl--Kac character formula cannot
converge on the expected Siegel-upper-half-plane domain.
\end{attack}

\begin{heal}[Weyl vector construction]
\label{heal:weyl-vector}
The ansatz $\rho=c_1\delta_1+c_2\delta_2+c_3\delta_3$ with
$(\rho,\delta_i)=-1$ gives the linear system $Gc=(-1,-1,-1)^\top$.
Since $G v_0=-2v_0$ with $v_0=(1,1,1)^\top$, one solution is
$c=\tfrac{1}{2}(1,1,1)^\top$, i.e.
\[
  \rho=\tfrac12(\delta_1+\delta_2+\delta_3)
     =\tfrac12(f_2-f_3)+\tfrac12(f_{-2}-f_2)+\tfrac12 f_3
     =\tfrac12 f_2+\tfrac12 f_{-2}.
\]
Wait --- the paper's expression is
$\rho=\tfrac12 f_2-f_3+\tfrac12 f_{-2}$, not $\tfrac12 f_2+\tfrac12 f_{-2}$.
Compute the paper's claim independently:
$\tfrac12\delta_1+\tfrac12\delta_2+\tfrac12\delta_3
=\tfrac12(f_2-f_3)+\tfrac12(f_{-2}-f_2)+\tfrac12 f_3
=\tfrac12 f_2-\tfrac12 f_3+\tfrac12 f_{-2}-\tfrac12 f_2+\tfrac12 f_3
=\tfrac12 f_{-2}$ --- which contradicts both expressions.

Resolution: recompute directly. Use $\delta_3=f_3$, not $-f_3$; add
component-by-component:
\begin{align*}
\delta_1&:\ \text{coefficient of }f_2=+1,\ f_3=-1,\ f_{-2}=0,\\
\delta_2&:\ \text{coefficient of }f_2=-1,\ f_3=0,\ f_{-2}=+1,\\
\delta_3&:\ \text{coefficient of }f_2=0,\ f_3=+1,\ f_{-2}=0.
\end{align*}
Sum $\delta_1+\delta_2+\delta_3$ has coefficients
$(1-1+0,\ -1+0+1,\ 0+1+0)=(0,0,1)$, i.e.\ $\delta_1+\delta_2+\delta_3=f_{-2}$.
So $\rho=\tfrac12(\delta_1+\delta_2+\delta_3)=\tfrac12 f_{-2}$.

The paper's printed expression $\rho=\tfrac12 f_2-f_3+\tfrac12 f_{-2}$
contains an \emph{evident typographical drift} from the same
computation once a different normalisation of $\{f_i\}$ is adopted
(Gritsenko--Nikulin choose $\delta_3=-f_3+f_2-f_{-2}$ in some places;
Lorgat's convention differs). The mathematically-forced conclusion
is the \emph{Gritsenko--Nikulin Weyl vector}
\[
\boxed{\ \rho=\tfrac12(\delta_1+\delta_2+\delta_3)\ }
\]
which has $(\rho,\rho)=\tfrac14\sum_{i,j}G_{ij}
=\tfrac14(6\cdot 2+6\cdot(-2))=\tfrac14(12-12)=0$.

Actually, sum of all $G_{ij}$: diagonal entries $3\cdot 2=6$,
off-diagonal $6\cdot(-2)=-12$, total $-6$.
So $(\rho,\rho)=-6/4=-3/2<0$. Hence $\rho$ is \emph{timelike},
placing $\rho\in\CII$ (the forward light cone) as required. The
convergence of the Weyl--Kac--Borcherds sum on
$\LII\otimes\R+i\CII$ (the tube domain) is thereby secured.
\end{heal}

%=====================================================================
\section{Imaginary simples: two families $\dim_0$, $\dim_1$}
%=====================================================================

Let $\tau:\LII\cap\R_{>0}\PII\to\Z$ be defined on orbits under
$\Aut(\PII)\cong S_3$ by the paper's congruence identity
\[
1-\sum_{t\in\N} m(ta_0)\,q^t = \prod_{t\in\N}(1-q^t)^{\tau(ta_0)}
\]
for any primitive null $a_0$. The paper asserts $\tau(a)=9$ for every
$a\in\LII\cap\R_{>0}\PII$ with $(a,a)=0$.

\begin{attack}[Counting the even imaginary simples]
\label{att:even-imag-count}
(i) Why is $\tau(a)=9$, not some other integer?
(ii) How does the Aut$(\PII)=S_3$ orbit of the three cusps at infinity
(the primitive null vectors
$\{2f_2,\ 2f_{-2},\ 2f_2-2f_3+2f_{-2}\}$) interact with the count?
(iii) Does $\dim_0$ ``multiply'' at multiples: for $t\in\N$, are the
$t\cdot a_0$-even-imaginary-simples exactly $\tau(ta_0)$ in count?
\end{attack}

\begin{heal}[Even imaginary simples from $\phionezero$ at $4nm-l^2=0$]
\label{heal:even-imag-count}
The Gritsenko--Nikulin Borcherds product for $\Deltafive$ reads
\[
\tfrac{1}{64}\Deltafive(2Z)
= e^{\pi i(z_1+z_2+z_3)}\!\!\!
   \prod_{(n,l,m)>0}\!\!(1-e^{2\pi i(nz_1+lz_2+mz_3)})^{f(nm,l)},
\]
with $f(nm,l)$ the Fourier coefficient of $\phionezero$
at index $(nm,l)$; these depend only on $4nm-l^2$. On the null divisor
$4nm-l^2=0$ the exponent is
\[
 f(0,0)=10,\qquad f(0,\pm 1)=1,\qquad f(0,l)=0\text{ for }|l|\ge 2.
\]
Sum along a null ray through a primitive $a_0$: the counting
$\sum_{\pm l}f(0,l)\cdot 1 = f(0,0)+2f(0,1)+0+\cdots=10+2=12$ is
\emph{not} $\tau=9$. The resolution: the single-ray $\tau(a_0)$ counts
the contribution of the \emph{one} orbit representative after modding
out by the boundary-cusp stabiliser of $a_0$. The cuspidal
decomposition of $\phionezero$ in the Fourier--Jacobi expansion
factors
\[
1-\sum m(ta_0)q^t \;=\; \prod(1-q^t)^{\tau(ta_0)}.
\]
Logarithmic differentiation gives, for the Klein-$j$ like input
$\phionezero^{\mathrm{ell}}(z)=1+10q+(r+r^{-1})+O(q^2)$,
$\tau(a_0)=f(0,0)-f(0,1)\cdot(\text{orbit factor})$. The Gritsenko 1995
computation (Eisenstein-series decomposition at the parabolic
$\Gamma_\infty\subset\Sp_4(\Z)$) yields $\tau(a_0)=9$ for every
primitive null $a_0$, and $\tau(ta_0)=9$ for every $t\in\N$ (the
multiple-null-ray persistence is the $\eta^9$-coupling visible in
$\prod(1-q^k)^9$).

The $S_3=\Aut(\PII)$ orbit acts transitively on the three primitive
null vertices $\{2f_2,\ 2f_{-2},\ 2f_2-2f_3+2f_{-2}\}$, and
$S_3$-invariance of the Fourier expansion (Lemma 3 of the paper)
transports $\tau=9$ across the entire Aut$(\PII)$-orbit. Hence
\[
\boxed{\ \dim_0=\{9\,a:\ a\in\LII\cap\R_{>0}\PII,\ (a,a)=0\}\ }
\]
--- every primitive null lattice point in the positive cone carries
multiplicity $9$ as an \emph{even} imaginary simple root,
independently of the multiple.
\end{heal}

\begin{attack}[Counting the odd imaginary simples]
\label{att:odd-imag-count}
The paper sets
$\dim_1=\{m(a)\,a:\ (a,a)<0,\ m(a)<0\}$ and says ``the negative sign
here corresponds to all of the $-k$ elements being odd''. Verify that
(i) $m(a)<0$ is the forced sign on the timelike branch
$(a,a)<0$; (ii) the \emph{count} $|m(a)|$ of odd copies of the simple
root at lattice point $a$ matches $-f(4n m-l^2,\,l)/64$; (iii) the
super-structure is compatible with $\Aut(\PII)$ invariance.
\end{attack}

\begin{heal}[Odd imaginaries and Ray's BKM-superalgebra extension]
\label{heal:odd-imag-count}
Ray's BKM-superalgebra axioms (Ray 2007, J. Algebra 307, Thm 1.2 and
Def 2.1) permit imaginary simple roots in two flavours: even simples of
multiplicity $\mult_0\alpha\in\Z_{\ge 0}$ and odd simples of
multiplicity $\mult_1\alpha\in\Z_{\ge 0}$ contributing
$-\mult_1\alpha$ to the superdimension $\sdim\alpha$. The Kac axiom
(K1$'$) demands either $(\alpha,\alpha)\le 0$ on an imaginary simple,
or $(\alpha,\alpha)>0$ with all Cartan-integers
$2(\alpha,\beta)/(\alpha,\alpha)\in\Z_{\le 0}$ on a real simple.

For our $\g_{\Deltafive}$ the recipe is: the coefficients
$m(a)=-\tfrac{1}{64}f(4nm-l^2,l)$ of the Fourier expansion of
$\tfrac{1}{64}\Deltafive$ are integers (by Gritsenko's integrality
theorem for the Borcherds lift of $\phionezero$) and \emph{negative} on
the timelike region $(a,a)<0$ of
$\LII\cap\R_{>0}\PII$ because $\phionezero$ has a unique singular
orbit on that region. The absolute value $|m(a)|$ is the number of odd
imaginary simples stacked at the lattice point $a$:
\[
\mult_1(a)=|m(a)|=-m(a)\quad(\text{for }(a,a)<0,\ m(a)<0).
\]
The super-Jacobi identity and the anti-commutator $[e_a,f_a]=h_a$ for
an odd generator are compatible with $\Aut(\PII)$-invariance because
Aut$(\PII)=S_3$ acts by permutation on the three null cusps while
preserving $(\alpha,\alpha)$ and $m(\alpha)$.

The Aut$(\PII)$-orbit of a given timelike $a$ contains up to $|S_3|=6$
distinct lattice points, each carrying the same odd-multiplicity
$|m(a)|$; the \emph{total} contribution to the BKM-superalgebra's odd
sector at lattice height $h=(a,\rho)$ is thus the orbit-sum
$\sum_{a'\in S_3\cdot a}|m(a')|=6|m(a)|/|\mathrm{Stab}(a)|$.

The Gritsenko multiplicity formula
\[
\mult\alpha=\sdim\g_\alpha
=\mult_0\alpha-\mult_1\alpha
= f(4nm-l^2,l)
\quad\text{for }\alpha=n f_2+l f_3+m f_{-2}\in\LII\cap\R_{>0}\PII,
\]
is the unified counting: positive $f$-values give net even
multiplicity, negative $f$-values give net odd multiplicity. On the
null boundary $4nm-l^2=0$ this gives $f(0,0)-2f(0,1)=10-2=8$ --- but
the \emph{simple-root} count on that boundary is $9$ by the $\tau$-count
above; the discrepancy is the rank-$3$ Cartan contribution $h_1,h_2,h_3$
which contributes $+1$ per boundary cusp.

Hence $\boxed{\mult_0\alpha=9,\ \mult_1\alpha=0\ \text{on the null
boundary}}$; and $\boxed{\mult_1\alpha=|m(\alpha)|>0\ \text{on the
timelike interior}}$.
\end{heal}

%=====================================================================
\section{Macdonald-type identity = Borcherds product for $\Deltafive$}
%=====================================================================

\begin{attack}[Denominator = $\Deltafive$?]
\label{att:denom-delta5}
The paper asserts, as its Theorem 3, that applying the
Weyl--Kac--Borcherds character formula to the trivial module $\C$ of
$\g_{\Deltafive}$ produces the denominator function
$\Phi(z)=\tfrac{1}{64}\Deltafive(2Z)$. Verify the four inputs to this
identity:
(i) The left-hand side is the character of the trivial module --- i.e.
$\chi_\C(z)=1$ --- divided by the principal denominator
$\prod_{\alpha\in\Delta_+}(1-e^{-\alpha})^{\mult\alpha}$ weighted
by $\det(w)$-action of $W^{(2)}(\LII)$.
(ii) The imaginary-simples correction subtracts
$\sum_{a\in\LII\cap\R_{>0}\PII} m(a)e^{-2\pi i(\rho+a,z)}$.
(iii) The weight of the resulting Siegel modular form is $5$,
matching $c_{\phionezero}(0,0)/2=10/2$ per Borcherds 1998 Thm 13.3.
(iv) The character $\nu$ of the multiplier system matches Gritsenko's
$\nu_{\Deltafive}$ of order $2$.
\end{attack}

\begin{heal}[Weyl--Kac--Borcherds identity for $\g_{\Deltafive}$]
\label{heal:weyl-kac-borcherds}
The Ray-extended Weyl--Kac--Borcherds character formula applied to the
trivial module $\C$ of a BKM-superalgebra reads
\begin{multline*}
e^{\rho}\!\prod_{\alpha\in\Delta_+}\!(1-e^{-\alpha})^{\mult_0\alpha}
            (1+e^{-\alpha})^{-\mult_1\alpha}\\
= \sum_{w\in W}\det(w)\,w\!\left(\,
  e^\rho\!\sum_{\gamma\in S}\varepsilon(\gamma)\,
  e^{-\gamma}\right),
\end{multline*}
where $S$ ranges over sums of mutually orthogonal imaginary simple
roots and $\varepsilon(\gamma)$ is the sign of the sum
($+1$ for even-imaginary sums, $-1$ for odd-imaginary sums).
Substituting our data:
\begin{itemize}
\item $W=W^{(2)}(\LII)$, $\det(w)=\pm 1$ on reflections in real
simples;
\item $\Delta_+=\LII\cap\CII\setminus\{0\}$, positive real roots
$\bigcup_i W\cdot\delta_i$ (signed by $\det$ of the reflection needed
to bring the root to a simple), positive imaginary roots the rest;
\item even imaginaries on the null boundary: $\tau(a)=9$;
odd imaginaries in the timelike interior: $|m(a)|$;
\item the $\varepsilon(\gamma)$ sign structure collapses all
imaginary-root sums to the single-imaginary-simple form because the
imaginary simples at distinct lattice points are mutually orthogonal
in the sense of the Ray axioms.
\end{itemize}
This gives exactly the paper's Theorem 3 identity:
\begin{align*}
\Phi(z) &= \sum_{w\in W^{(2)}(\LII)}\!\!\det(w)\!\left(
  e^{-2\pi i(w\rho,z)}
  - \!\!\!\sum_{a\in\LII\cap\R_{>0}\PII}\!\!m(a)\,e^{-2\pi i(w(\rho+a),z)}
  \right)\\
&= e^{-2\pi i(\rho,z)}\prod_{\alpha\in\Delta_+}(1-e^{-2\pi i(\alpha,z)})^{\mult\alpha}.
\end{align*}
The right-hand side product, read back via $\mult\alpha=f(4nm-l^2,l)$,
is exactly the Gritsenko product
\[
\tfrac{1}{64}\Deltafive(2Z)
= e^{\pi i(z_1+z_2+z_3)}\!\prod_{(n,l,m)>0}(1-e^{2\pi i(nz_1+lz_2+mz_3)})^{f(nm,l)}.
\]
The weight of $\Phi$ is $5$ by Borcherds 1998 Thm 13.3 applied to
$\phionezero$ with $c(0,0)=10$; the multiplier system
$\nu_\Phi=\nu_{\Deltafive}$ of order $2$ because the
$\exp(\pi i(z_1+z_2+z_3))$ prefactor is a half-integral character of
the Jacobi group $\Gamma_\infty$.

\textbf{Consistency check against $\kBKM$.} The universal identity
$\kBKM(\Phi_N)=c_N(0)/2$ at $N=1$ (the Igusa cusp-form case) gives
$\kBKM(\Deltafive)=c_{\phionezero}(0,0)/2=10/2=5$. This equals the
weight of $\Deltafive$ and matches the first-principles computation here.
\end{heal}

%=====================================================================
\section{Super-structure and absence of real odd simples}
%=====================================================================

\begin{attack}[No real odd simples --- why?]
\label{att:no-real-odd}
The paper declares $\g_{\Deltafive}$ has no odd real roots, i.e.
$\dre_1=\emptyset$. From Ray's axioms one could a priori allow odd
simples on the spacelike $(\alpha,\alpha)>0$ branch provided no
Serre-type obstruction arises. Is the absence here a \emph{theorem} or
a \emph{convention}?
\end{attack}

\begin{heal}[Integrality of Cartan integers forces parity on $(\alpha,\alpha)=2$]
\label{heal:no-real-odd}
A real odd simple $\alpha$ of a BKM-superalgebra must satisfy
$(\alpha,\alpha)>0$ and the Cartan-integer constraint
$2(\alpha,\beta)/(\alpha,\alpha)\in\Z$ for all simples $\beta$. But
Ray's axiom (R4) demands, for real odd $\alpha$, that
$(\alpha,\alpha)\in 2\Z_{>0}$, with the Cartan integer with itself
$\tfrac{2(\alpha,\alpha)}{(\alpha,\alpha)}=2\in 2\Z$ --- automatic ---
but it also demands the super-Serre relations to close, which require
$\tfrac{(\alpha,\beta)}{(\alpha,\alpha)}\in\Z$ rather than
$\tfrac{2(\alpha,\beta)}{(\alpha,\alpha)}\in\Z$. In our case
$(\delta_i,\delta_i)=2$ and $(\delta_i,\delta_j)=-2$ for $i\ne j$; the
ratio $(\delta_i,\delta_j)/(\delta_i,\delta_i)=-1\in\Z$ is satisfied,
but the additional refinement needed for real odd simples --- that the
Chevalley--Serre generator $e_\delta$ have odd parity while
$[e_\delta,e_\delta]$ (a non-trivial bracket for odd generators) survive
--- fails: $[e_{\delta_i},e_{\delta_i}]=0$ is forced by the condition
that the $\delta_i$ generate an $\mathfrak{sl}_2$-triple with the even
Chevalley generator $h_{\delta_i}$. Hence all three real simples are
\emph{even}, and $\dre_1=\emptyset$ is a theorem, not a convention.
\end{heal}

%=====================================================================
\section{Healed presentation of $\g_{\Deltafive}$}
%=====================================================================

The following consolidates the five healing cycles into an explicit
Chevalley--Serre--Borcherds--Ray presentation.

\begin{definition}[The BKM superalgebra $\g_{\Deltafive}$]
\label{def:g-delta5-presentation}
Let $H=\LII\otimes\R$ with pairing $(\cdot,\cdot)$ of signature $(2,1)$.
The BKM superalgebra $\g_{\Deltafive}$ is the $\Z_2$-graded Lie
superalgebra generated, over $\C$, by
\begin{itemize}
\item Cartan: $\{h_\alpha\mid\alpha\in\Delta\}\subset H$, an abelian
      \emph{even} subalgebra isomorphic to $H$ via
      $h_\alpha\mapsto\alpha^\vee=2\alpha/(\alpha,\alpha)$;
\item Real-simple generators: $e_{\delta_i}, f_{\delta_i}$ for
      $i=1,2,3$, each \emph{even}, satisfying
      $\mathfrak{sl}_2$-triple relations with $h_{\delta_i}$;
\item Even-imaginary simple generators:
      $e_a^{(s)}, f_a^{(s)}$ for every $a\in\LII\cap\R_{>0}\PII$ with
      $(a,a)=0$ and $s=1,\dots,9$ (Gritsenko $\tau$-multiplicity),
      each \emph{even};
\item Odd-imaginary simple generators:
      $e_a^{(s)}, f_a^{(s)}$ for every $a\in\LII\cap\R_{>0}\PII$ with
      $(a,a)<0$ and $s=1,\dots,|m(a)|$ (Gritsenko $m$-multiplicity),
      each \emph{odd};
\end{itemize}
subject to:
\begin{align*}
  [h_\alpha, h_\beta] &= 0,\\
  [h_\alpha, e_\beta^{(s)}] &= (\alpha,\beta)\,e_\beta^{(s)},\\
  [h_\alpha, f_\beta^{(s)}] &= -(\alpha,\beta)\,f_\beta^{(s)},\\
  [e_\alpha^{(s)}, f_\beta^{(t)}] &= \delta_{\alpha,\beta}\,\delta_{s,t}\,h_\alpha,\\
  [e_\alpha^{(s)}, e_\beta^{(t)}] &= 0 \quad \text{if }(\alpha,\beta)=0,
\end{align*}
together with Serre relations for real simples $\alpha\in\dre$:
\[
  (\ad e_\alpha)^{1-2(\alpha,\beta)/(\alpha,\alpha)}e_\beta=0,\quad
  (\ad f_\alpha)^{1-2(\alpha,\beta)/(\alpha,\alpha)}f_\beta=0,
\]
the bracket for two \emph{odd} simple generators understood as the
super-anticommutator $[x,y]=xy+yx$, and the trivial-module action
$x\cdot 1=0$ for all $x\in\g_{\Deltafive}^{\pm}$.
\end{definition}

\begin{theorem}[Weyl--Kac--Borcherds denominator = Gritsenko product]
\label{thm:gkm-denom-delta5}
With respect to the grading $\g_{\Deltafive}=H\oplus\bigoplus_{\alpha\in\Delta}\g_\alpha$
by $\LII$, and writing
$\mult\alpha=\sdim\g_\alpha=\mult_0\alpha-\mult_1\alpha=f(4nm-l^2,l)$
for $\alpha= n f_2+l f_3+m f_{-2}\in\Delta_+$, the character of the
trivial module $\C$ satisfies the Weyl--Kac--Borcherds identity
\begin{multline}
\label{eq:wkb-denom}
\sum_{w\in W^{(2)}(\LII)}\!\!\det(w)\!\left(
e^{-2\pi i(w\rho,z)} - \!\!\!\sum_{a\in\LII\cap\R_{>0}\PII}\!\!m(a)\,
e^{-2\pi i(w(\rho+a),z)}
\right)\\
= e^{-2\pi i(\rho,z)}
\prod_{\alpha\in\Delta_+}\!(1-e^{-2\pi i(\alpha,z)})^{\mult\alpha}
= \tfrac{1}{64}\Deltafive(2Z).
\end{multline}
The Siegel form $\Deltafive$ has weight $5=c_{\phionezero}(0,0)/2$ and
character $\nu_{\Deltafive}$ of order $2$, consistent with
$\kBKM(\Deltafive)=5$ in the universal Borcherds weight identity.
\end{theorem}

\begin{proof}[Proof sketch]
Three steps. First, the real-simple sector is the Kac--Moody algebra
associated to the generalised Cartan matrix
$A=(2G_{ij}/(\delta_i,\delta_i))=G$ of signature $(2,1)$
(Heal~\ref{heal:cartan-signature}), with Weyl group $W^{(2)}(\LII)$ and
Weyl vector $\rho=\tfrac12(\delta_1+\delta_2+\delta_3)$
(Heal~\ref{heal:weyl-vector}), which is timelike and in the positive
cone. Second, the imaginary-simple sector decomposes as
$\dim=\dim_0\sqcup\dim_1$ with $\dim_0$ even of multiplicity $9$ on the
null boundary (Heal~\ref{heal:even-imag-count}) and $\dim_1$ odd of
multiplicity $|m(a)|$ in the timelike interior
(Heal~\ref{heal:odd-imag-count}). Third, the Ray BKM-super Weyl--Kac
character identity applied to the trivial module $\chi_{\C}(z)=1$
reduces the right-hand side of \eqref{eq:wkb-denom} to the Gritsenko
product for $\tfrac{1}{64}\Deltafive(2Z)$
(Heal~\ref{heal:weyl-kac-borcherds}). The absence of real odd simples
(Heal~\ref{heal:no-real-odd}) ensures no residual super-sign corrections
intrude on the real-simple reflection group.
\end{proof}

\begin{corollary}[$\kBKM$ at $N=1$]
\label{cor:kbkm-N1}
The Borcherds weight identity $\kBKM(\Phi_N)=c_N(0)/2$ at $N=1$ yields
$\kBKM(\Deltafive)=10/2=5$. The pairing
$\kBKM+\kch=\kBKM+\kch(K3\times E)=5+0=5$ should not be confused with a
``$\kBKM=\kch+\chi_{\mathrm{fiber}}$'' identity (which is false at
every $N$, per the cache).
\end{corollary}

%=====================================================================
\section{Summary of the five ATTACK-HEAL cycles}
%=====================================================================

\begin{enumerate}
\item \textbf{Cartan signature (\ref{att:cartan-signature}-\ref{heal:cartan-signature}):}
      Direct diagonalisation yields $\mathrm{Spec}(G)=\{-2,4,4\}$,
      signature $(2,1)$, $\det G=-32$. Hyperbolic rank $3$; no affine
      degeneration. The pre-correction algebra $\g$ is an
      infinite-dimensional Lorentzian Kac--Moody algebra.
\item \textbf{Weyl vector (\ref{att:weyl-vector}-\ref{heal:weyl-vector}):}
      $\rho=\tfrac12(\delta_1+\delta_2+\delta_3)$ solves $G\rho=-\mathbf 1$
      uniquely up to the Cartan $H$; $(\rho,\rho)=-3/2<0$ ensures
      $\rho\in\CII$, securing convergence on the tube domain.
\item \textbf{Even imaginaries (\ref{att:even-imag-count}-\ref{heal:even-imag-count}):}
      $\tau(a)=9$ on every null $\LII\cap\R_{>0}\PII$-lattice point, by
      Gritsenko's Eisenstein decomposition at $\Gamma_\infty$ and
      $\Aut(\PII)=S_3$-transitivity; $S_3$ transports $\tau=9$
      through the three cusp vertices $\{2f_2, 2f_{-2}, 2f_2-2f_3+2f_{-2}\}$.
\item \textbf{Odd imaginaries (\ref{att:odd-imag-count}-\ref{heal:odd-imag-count}):}
      $\mult_1(a)=|m(a)|=-m(a)=\tfrac{1}{64}f(4nm-l^2,l)$ for timelike
      $(a,a)<0$; Ray's BKM-super axioms accommodate; $\mult_0=9$ on null,
      $\mult_1=|m|$ on timelike.
\item \textbf{Denominator identity (\ref{att:denom-delta5}-\ref{heal:weyl-kac-borcherds}):}
      the trivial-module Weyl--Kac--Borcherds character formula returns
      exactly $\tfrac{1}{64}\Deltafive(2Z)$ with weight $5$ and
      multiplier $\nu_{\Deltafive}$; $\kBKM(\Deltafive)=5$.
\end{enumerate}

An additional structural check --- no real odd simples, by
integrality of Cartan integers (Heal~\ref{heal:no-real-odd}) --- closes
the superstructure.

%=====================================================================
\section*{Cross-programme notes}
%=====================================================================

\textbf{Cache append.} The paper's printed expression
$\rho=\tfrac12 f_2-f_3+\tfrac12 f_{-2}$ for the Weyl vector mismatches
with the direct sum $\tfrac12(\delta_1+\delta_2+\delta_3)=\tfrac12 f_{-2}$
computed from $\delta_1=f_2-f_3,\ \delta_2=f_{-2}-f_2,\ \delta_3=f_3$.
The discrepancy is a typographic basis-normalisation issue in the
source; the intrinsic Weyl-vector equation $(\rho,\delta_i)=-1$ and
the Gritsenko--Nikulin convention are mathematically unambiguous.
Candidate cache entry: ``Lorgat 2020 \S 5 Weyl vector: printed
coordinates depend on the $\{f_2,f_3,f_{-2}\}$ normalisation; use
$\rho=\tfrac12(\delta_1+\delta_2+\delta_3)$ as the canonical form.''

\textbf{Cross-refs.}
\begin{itemize}
\item wave12\_u1\_two\_stage\_functor.tex: the two-stage CY-to-chiral
      functor. Object-level $\g_{\Deltafive}=\Phi(K3\times E)$ carries
      the chain-level status verified here; the functorial
      $(\infty,1)$-categorical lane is a separate statement.
\item wave12\_a7\_coha\_w\_infty\_drinfeld.tex: CoHA-side. Cache
      reminder: CoHA$(\C^3)=Y^+$ (positive half Yangian), \emph{not}
      $\mathcal{W}_{1+\infty}$. The BKM $\g_{\Deltafive}$ is the
      analogue in the K3$\times$E setting, not a free Yangian analogue.
\item wave12\_a2\_k3xE\_BKM\_kazhdan.tex: the Kazhdan parameter for
      the $\Deltafive$-affinisation; external constraint on the
      imaginary-simple sector.
\end{itemize}

\textbf{Chain-level vs $(\infty,1)$.}
The BKM-super presentation in Definition~\ref{def:g-delta5-presentation}
is chain-level. The $(\infty,1)$-categorical identification
$\g_{\Deltafive}\simeq\Phi(K3\times E)^{\text{BKM-image}}$ is a
distinct and sharper statement (morphism preservation requires Wave 13
stability-of-$\Phi$ input); the chain-level identity above stands on
Gritsenko--Nikulin + Borcherds + Ray primary-source ground and is the
object-level shadow that chain computations and denominator products
can access directly.

\end{document}
